% =====================================================================
%  CAPITOLO 1 - INTRODUZIONE
%
%  NOTE PER LA COMPILAZIONE:
%  - I comandi \cite{...} rimandano a voci bibliografiche da inserire e
%    VERIFICARE in bibliografia.bib (vedi elenco in fondo al file).
%  - La figura dell'architettura va realizzata in PowerPoint/strumento
%    vettoriale ed esportata in PDF (linea guida relatore), poi inserita
%    al posto del placeholder indicato.
% =====================================================================

\chapter{Introduzione}
\label{cap:introduzione}

\section{Contesto e motivazione}
\label{sec:contesto}

L'obesit\`a infantile \`e riconosciuta come una delle pi\`u rilevanti sfide di
salute pubblica a livello globale e nazionale \cite{who_obesity}. In Italia, i
sistemi di sorveglianza dedicati alla popolazione in et\`a scolare rilevano una
quota significativa di bambini in condizione di eccesso ponderale, accompagnata
da una diffusa sottostima del fenomeno da parte delle famiglie
\cite{okkio}. La condizione, oltre alle conseguenze cliniche dirette, tende a
persistere nell'et\`a adulta e si associa a un aumento del rischio di patologie
croniche, rendendo la prevenzione precoce un obiettivo prioritario delle linee
guida nazionali e internazionali \cite{linee_guida_ministero}.

La prevenzione dell'obesit\`a infantile non \`e responsabilit\`a di un singolo
attore: coinvolge i genitori nella gestione delle abitudini alimentari
quotidiane, gli insegnanti nel contesto educativo e della ristorazione
scolastica, e i professionisti sanitari nella valutazione e nella presa in
carico. Ciascuno di questi interlocutori ha esigenze informative diverse per
linguaggio, livello di dettaglio tecnico e tipo di contenuto. L'accesso a
informazioni affidabili e contestualizzate rappresenta dunque un fattore
abilitante per interventi di prevenzione efficaci.

\section{Il problema}
\label{sec:problema}

Negli ultimi anni i modelli linguistici di grandi dimensioni
(\emph{Large Language Models}, LLM) hanno reso possibile l'interazione in
linguaggio naturale con grandi basi di conoscenza, offrendo un canale di
accesso immediato e conversazionale all'informazione. Tuttavia, l'impiego
diretto di un LLM generalista in un dominio sanitario presenta criticit\`a non
trascurabili. In primo luogo, i modelli generativi possono produrre affermazioni
plausibili ma prive di fondamento nelle fonti, fenomeno noto come
\emph{allucinazione} \cite{ji_hallucination}: in un contesto in cui le risposte
riguardano la salute di un minore, un'informazione errata o inventata pu\`o avere
conseguenze concrete. In secondo luogo, un sistema rivolto al pubblico generale
non deve sostituirsi al professionista sanitario, e va quindi impedito che
formuli diagnosi, prescriva terapie o restituisca contenuti clinicamente
inappropriati.

Il paradigma \emph{Retrieval-Augmented Generation} (RAG) \cite{lewis_rag}
affronta il primo problema vincolando la generazione a un insieme di documenti
recuperati dinamicamente da una base di conoscenza verificata, riducendo lo
spazio in cui il modello pu\`o produrre contenuti non fondati. Il secondo
problema --- la sicurezza e l'appropriatezza clinica delle risposte --- richiede
invece meccanismi di controllo dedicati, comunemente indicati come
\emph{guardrail}, che intervengano sia sulle richieste in ingresso sia sulle
risposte in uscita. La combinazione di questi due aspetti, unita alla necessit\`a
di adattare la comunicazione a interlocutori eterogenei, definisce il problema
affrontato in questa tesi.

\section{Obiettivi}
\label{sec:obiettivi}

L'obiettivo del presente lavoro \`e la progettazione e la realizzazione di un
agente conversazionale in lingua italiana che fornisca supporto informativo
sulla prevenzione dell'obesit\`a infantile (fascia 0--10 anni) a una pluralit\`a
di interlocutori, garantendo al contempo l'aderenza alle fonti e la sicurezza
dei contenuti. Tale obiettivo generale si articola nei seguenti obiettivi
specifici:

\begin{itemize}
    \item realizzare una pipeline RAG che recuperi le informazioni da un corpus
    di linee guida e documenti scientifici italiani e che fondi le risposte
    esclusivamente su tale corpus, con citazione esplicita delle fonti;
    \item adattare automaticamente registro, lessico e contenuto delle risposte
    al profilo dell'interlocutore (genitore, insegnante, pediatra,
    nutrizionista, ricercatore);
    \item integrare la pipeline in un'architettura di sicurezza che prevenga la
    generazione di contenuti clinicamente inappropriati e l'uso improprio del
    sistema;
    \item valutare quantitativamente la qualit\`a delle risposte, la robustezza
    dei meccanismi di sicurezza e l'efficacia dell'adattamento al profilo.
\end{itemize}

\section{Contributi}
\label{sec:contributi}

I principali contributi di questa tesi sono i seguenti:

\begin{itemize}
    \item un agente conversazionale RAG in lingua italiana per la prevenzione
    dell'obesit\`a infantile, costruito su un corpus curato di documenti
    istituzionali e scientifici e dotato di un meccanismo di recupero ibrido
    che combina ricerca semantica vettoriale e ricerca lessicale;
    \item un meccanismo di adattamento al profilo dell'interlocutore, che
    differenzia la comunicazione su cinque categorie di stakeholder;
    \item un'architettura di \emph{guardrail} a difesa in profondit\`a, composta
    da controlli sull'input, verifica di ancoraggio alle fonti a tempo di
    esecuzione e controlli sull'output, con registrazione tracciabile degli
    interventi;
    \item una valutazione sperimentale multi-prospettica che misura la fedelt\`a
    e la pertinenza delle risposte (tramite la metodologia RAGAS), la robustezza
    del livello di sicurezza (tramite una campagna di \emph{red-teaming}) e
    l'efficacia dell'adattamento al profilo.
\end{itemize}

L'architettura complessiva del sistema, dall'ingresso della richiesta dell'utente
fino alla restituzione della risposta, \`e illustrata in
Figura~\ref{fig:architettura}. Come si osserva nella figura, la richiesta
attraversa in sequenza il livello di controllo in ingresso, la pipeline RAG vera
e propria e il livello di controllo in uscita, secondo il principio della difesa
in profondit\`a discusso nei capitoli successivi.

% ---------------------------------------------------------------------
% FIGURA ARCHITETTURA
% Da realizzare in PowerPoint/strumento vettoriale ed esportare in PDF
% (linea guida relatore: niente immagini raster sgranate).
% Salvare il file come tesi/immagini/architettura.pdf e decommentare.
% ---------------------------------------------------------------------
% \begin{figure}[htbp]
%     \centering
%     \includegraphics[width=0.95\textwidth]{immagini/architettura.pdf}
%     \caption{Architettura complessiva del sistema: la richiesta dell'utente
%     attraversa il livello di guardrail in ingresso, la pipeline RAG
%     (rilevamento del profilo, riformulazione della query, recupero ibrido e
%     generazione) e il livello di guardrail in uscita.}
%     \label{fig:architettura}
% \end{figure}

\section{Struttura della tesi}
\label{sec:struttura}

Il resto della tesi \`e organizzato come segue. Il
Capitolo~\ref{cap:background} introduce i fondamenti teorici e lo stato
dell'arte rilevanti per il lavoro: i modelli linguistici di grandi dimensioni,
il paradigma RAG, i meccanismi di guardrail e il dominio applicativo della
prevenzione dell'obesit\`a infantile. Il Capitolo~\ref{cap:progettazione}
descrive la progettazione dell'architettura del sistema e le scelte
metodologiche adottate. Il Capitolo~\ref{cap:implementazione} dettaglia
l'implementazione dei singoli componenti. Il Capitolo~\ref{cap:valutazione}
presenta la metodologia e i risultati della valutazione sperimentale. Infine,
il Capitolo~\ref{cap:conclusioni} riassume i risultati ottenuti, ne discute i
limiti e delinea i possibili sviluppi futuri.

% =====================================================================
%  RIFERIMENTI BIBLIOGRAFICI CITATI IN QUESTO CAPITOLO (da inserire e
%  VERIFICARE in bibliografia.bib):
%    who_obesity            -> dati/definizioni OMS sull'obesita' infantile
%    okkio                  -> sistema di sorveglianza OKkio alla SALUTE
%    linee_guida_ministero  -> Linee di indirizzo Ministero della Salute
%    ji_hallucination       -> survey sulle allucinazioni nei modelli generativi
%    lewis_rag              -> Lewis et al., Retrieval-Augmented Generation, 2020
%
%  ETICHETTE DI CAPITOLO ATTESE (da definire con \label nei rispettivi file):
%    cap:background, cap:progettazione, cap:implementazione,
%    cap:valutazione, cap:conclusioni
% =====================================================================
